% authority: science_draft
\documentclass[11pt]{article}

\usepackage[margin=1in]{geometry}
\usepackage{amsmath,amssymb,amsthm}
\usepackage{booktabs}
\usepackage{enumitem}
\usepackage{longtable}

\newtheorem{definition}{Definition}
\newtheorem{theorem}{Theorem}
\newtheorem{proposition}{Proposition}
\newtheorem{remark}{Remark}
\newtheorem{nonconclusion}{Non-conclusion}

\newcommand{\Ksrc}{K_{\mathrm{src}}}
\newcommand{\Conf}{\mathcal C_{\mathrm{mat}}}
\newcommand{\Trans}{\mathcal T_{\mathrm{mat}}}
\newcommand{\Control}{\mathfrak C}
\newcommand{\Enabled}{\mathcal E}
\newcommand{\Book}{\mathcal B}
\newcommand{\Sect}{\mathrm{Sec}_{\mathrm{src}}}
\newcommand{\Kernel}{K_{\Control}}
\newcommand{\Update}{U_{\Control}}
\newcommand{\Malformed}{\bot_K}
\newcommand{\Undetermined}{\mathsf{undetermined}_K}

\title{P7-T02 Source-Matter Finite Transition Kernel Candidate v1}
\author{Candidate Constructor Packet RT-20260728-001}
\date{2026-07-28}

\begin{document}
\maketitle

\section{Control status and decisive result}

This artifact is the fused science-draft output of
\texttt{AJ-RT-20260728-001-001}. It constructs
\[
  \mathrm{SourceMatterFiniteTransitionKernelCandidate}_{v1}
\]
as a \emph{draft/control, proposal-only} source-transition-law candidate.
The candidate supplies a finite enabled-transition system, a dimensionless
row-normalized bookkeeping kernel, a reproducible finite update equation, an
analytic positive control, exact normalization and closure, an inherited
sector-support theorem, source-presentation equivariance, independent and
interacting controls, and exact null, degenerate, malformed, and
underdetermined branches.

The candidate is not derived by the current ontology and is not adopted.
Adoption remains
\texttt{blocked\_adoption\_open\_continuation}. Kernel coefficients are not
physical probabilities, rates, frequencies, or empirical propensities. The
composition index is not physical time, duration, or causal order. The
artifact does not establish physical matter, a physical conservation law,
detector semantics, effective geometry, stress energy, a matter action,
universal coupling, Einstein equations, exact-GR recovery, benchmark
promotion, proof authority, publication, completed derivation, a global
no-go theorem, or future source-extension impossibility.

\section{Fixed P7-T01 basis}

The construction holds
\(\mathrm{SourceMatterIncidenceCandidate}_{v1}\) fixed. Its admissible
configurations are \(\Psi=(m,z)\), with occupation multiplicities \(m\),
closed defect chain \(z\), derived charge chain \(q(\Psi)\), and sector
\[
  \Sect(\Psi)=\bigl(Q(\Psi),D(\Psi)\bigr),
  \qquad
  Q(\Psi)=\epsilon(q(\Psi)),
  \qquad
  D(\Psi)=[z].
\]
An admissible P7-T01 transition record is
\[
  T=(\Psi^-,J,w,\Pi,\Psi^+)
\]
with
\begin{equation}
  q(\Psi^+)-q(\Psi^-)=\partial_1J,
  \qquad
  z^+-z^-=\partial_2w,
  \label{eq:p7t01-constraints}
\end{equation}
and source-only update record \(\Pi\). Because
\(\epsilon\partial_1=0\) and boundaries do not change homology classes,
every such transition has
\[
  \Sect(\Psi^+)=\Sect(\Psi^-).
\]
This is inherited kinematic admissibility. It does not say which transition
is enabled, how alternatives are weighted, whether any event occurs, or how
updates compose. It is not a physical continuity equation or Noether law.

\section{Finite control systems}

\begin{definition}[Finite source-transition control system]
A finite control system is
\[
  \Control=(\Omega,\Enabled,s,t),
\]
where \(\Omega\) is a finite nonempty subset of admissible P7-T01
configurations, \(\Enabled\) is a finite set of admissible P7-T01 transition
records with source and target maps
\[
  s,t:\Enabled\longrightarrow\Omega,
\]
and every enabled row
\[
  \Enabled_x=\{T\in\Enabled:s(T)=x\}
\]
contains the typed identity record
\[
  \mathbf 1_x=(x,0,0,\Pi_{\mathrm{id}},x).
\]
Here \(\Pi_{\mathrm{id}}\) is the declared source-only record that leaves
occupations and labels unchanged. Thus
\[
  d_{\Control}(x)=|\Enabled_x|\geq 1.
\]
The finite set \(\Omega\), enabled rows, and identity convention are
proposal inputs. The identity record is a totality convention for the
bookkeeping update, not physical persistence, inertia, or a zero-duration
event. The finite control bound is not a physical cutoff, spatial region, or
finiteness claim about nature.
\end{definition}

\begin{definition}[Uniform finite transition kernel]
For a finite control system define
\begin{equation}
  \Kernel(T\mid x)=
  \begin{cases}
    d_{\Control}(x)^{-1},&T\in\Enabled_x,\\
    0,&T\notin\Enabled_x.
  \end{cases}
  \label{eq:kernel}
\end{equation}
The codomain is \(\mathbb Q_{\geq0}\). Uniformity is a minimal
presentation-neutral proposal: a row carries no ordering or other declared
source datum that could distinguish its members. There are no tunable
coupling constants and no dimensional scales. Row cardinalities and the
finite control domain are exact task-local data.
\end{definition}

The support relation is primary. Equation~\eqref{eq:kernel} is one
proposal-only completion of that relation, not a claim that alternatives
occur randomly or with empirical frequency.

\section{Reproducible finite update equation}

\begin{definition}[Bookkeeping state and update]
Let
\[
  \Book(\Omega)=\mathbb Q_{\geq0}^{\Omega},
  \qquad
  \|\mu\|_1=\sum_{x\in\Omega}\mu(x).
\]
Elements of \(\Book(\Omega)\) are finite rational bookkeeping weights, not
empirical ensembles. Define
\begin{equation}
  (\Update\mu)(y)
  =
  \sum_{x\in\Omega}\mu(x)
  \sum_{\substack{T\in\Enabled_x\\t(T)=y}}
  \Kernel(T\mid x).
  \label{eq:update}
\end{equation}
The recurrence
\[
  \mu_{n+1}=\Update\mu_n,\qquad n\in\mathbb N_0,
\]
counts formal compositions of the candidate update. It assigns no physical
time, duration, causal ordering, clock reading, or trajectory.
\end{definition}

\begin{theorem}[Finite normalization, closure, and inherited-sector support]
\label{thm:normalization}
For every finite control system \(\Control\):
\begin{enumerate}[leftmargin=*]
\item each row is normalized,
  \[
    \sum_{T\in\Enabled_x}\Kernel(T\mid x)=1;
  \]
\item \(\Update\) maps \(\Book(\Omega)\) to itself;
\item total bookkeeping mass is preserved,
  \[
    \|\Update\mu\|_1=\|\mu\|_1;
  \]
\item the unit-mass subset
  \(\Book_1(\Omega)=\{\mu:\|\mu\|_1=1\}\) is closed under
  \(\Update\);
\item if \(\mu\) is supported in one fixed P7-T01 sector \(\sigma\), then
  \(\Update\mu\) is supported in \(\sigma\).
\end{enumerate}
\end{theorem}

\begin{proof}
The first statement is
\[
  \sum_{T\in\Enabled_x}\frac{1}{d_{\Control}(x)}
  =
  \frac{|\Enabled_x|}{d_{\Control}(x)}=1.
\]
Every sum in~\eqref{eq:update} is finite and nonnegative rational, proving
closure. Reordering the finite sums gives
\[
\begin{aligned}
  \|\Update\mu\|_1
  &=
  \sum_{y\in\Omega}\sum_{x\in\Omega}\mu(x)
  \sum_{\substack{T\in\Enabled_x\\t(T)=y}}
  \Kernel(T\mid x)\\
  &=
  \sum_{x\in\Omega}\mu(x)
  \sum_{T\in\Enabled_x}\Kernel(T\mid x)
  =
  \sum_{x\in\Omega}\mu(x).
\end{aligned}
\]
The unit-mass claim follows. Finally, every transition in the support obeys
the inherited constraints~\eqref{eq:p7t01-constraints}, hence has equal
sector coordinates at source and target. Therefore no update term can leave
the sector supporting \(\mu\).
\end{proof}

\begin{remark}[New and inherited content]
Row normalization, finite closure, total bookkeeping-mass preservation, and
the update equation are new consequences of the candidate kernel. Equality
of \(Q\) and \(D\) across each supported edge is inherited from P7-T01.
The distribution-level sector-support statement combines those layers; it is
not an independent physical conservation law.
\end{remark}

\section{Source-presentation equivariance}

\begin{definition}[Transported control system]
Let \(h\) be a compatible P7-T01 source-presentation isomorphism. It transports
configurations and transition records. A control system \(\Control'\) is the
transport of \(\Control\) when \(h:\Omega\to\Omega'\) and
\(h:\Enabled\to\Enabled'\) are bijections satisfying
\[
  s'(hT)=h(sT),\qquad t'(hT)=h(tT),
\]
and \(h(\Enabled_x)=\Enabled'_{hx}\).
For bookkeeping states define
\[
  (h_*\mu)(hx)=\mu(x).
\]
\end{definition}

\begin{theorem}[Source-presentation equivariance]
\label{thm:equivariance}
For transported finite control systems,
\[
  d_{\Control'}(hx)=d_{\Control}(x),
  \qquad
  K_{\Control'}(hT\mid hx)=K_{\Control}(T\mid x),
\]
and
\begin{equation}
  h_*\Update=U_{\Control'}h_*.
  \label{eq:intertwining}
\end{equation}
Consequently, if a declared presentation automorphism preserves
\(\Control\), its invariant bookkeeping subspace is preserved by
\(\Update\).
\end{theorem}

\begin{proof}
The row bijection preserves cardinality, which proves the kernel identity.
Substituting \(x'=hx\), \(T'=hT\), and \(y'=hy\) into the finite sum
in~\eqref{eq:update} proves~\eqref{eq:intertwining}. If \(h_*\mu=\mu\),
then \(h_*\Update\mu=\Update h_*\mu=\Update\mu\).
\end{proof}

This theorem concerns candidate presentation changes. It is not target
covariance, diffeomorphism invariance, physical gauge symmetry, or general
\(\mathrm{EqSrc}\).

\section{Free positive control and analytic solution}

Take two source vertices \(u,v\), one link \(e\) with
\(\partial_1e=[v]-[u]\), coefficient group \(A_{\mathrm{int}}=\mathbb Z\),
one representation label \(\rho\) with \(\chi(\rho)=1\), and zero defect.
Let \(x_u\) have one occupation at \(u\) and \(x_v\) one at \(v\). Both have
sector \((1,[0])\). The nonidentity transition
\[
  T_e=(x_u,e,0,\Pi_e,x_v)
\]
is admissible because its continuity residual is zero.

Choose
\[
  \Omega=\{x_u,x_v\},\qquad
  \Enabled_{x_u}=\{\mathbf1_{x_u},T_e\},\qquad
  \Enabled_{x_v}=\{\mathbf1_{x_v}\}.
\]
Then
\[
  K(\mathbf1_{x_u}\mid x_u)
  =
  K(T_e\mid x_u)=\frac12,
  \qquad
  K(\mathbf1_{x_v}\mid x_v)=1.
\]
For \(\mu_0=\delta_{x_u}\),
\[
  \mu_1=\frac12\delta_{x_u}+\frac12\delta_{x_v},
\]
and the exact solution is
\begin{equation}
  \mu_n(x_u)=2^{-n},
  \qquad
  \mu_n(x_v)=1-2^{-n}.
  \label{eq:positive-solution}
\end{equation}
Indeed,~\eqref{eq:positive-solution} satisfies the recurrence because
\(\mu_{n+1}(x_u)=\mu_n(x_u)/2\) and
\(\mu_{n+1}(x_v)=\mu_n(x_v)+\mu_n(x_u)/2\).
All coefficients stay nonnegative and sum to one. This is a stable bounded
solution of a finite bookkeeping recurrence, not particle propagation,
physical relaxation, or a trajectory in time.

\section{Independent and interacting branches}

\begin{proposition}[Independent product factorization]
For finite control systems \(\Control_1,\Control_2\), define the independent
product by
\[
  \Omega_{12}=\Omega_1\times\Omega_2,\qquad
  \Enabled_{(x_1,x_2)}
  =
  \Enabled^1_{x_1}\times\Enabled^2_{x_2}.
\]
Then
\[
  K_{12}(T_1,T_2\mid x_1,x_2)
  =
  K_1(T_1\mid x_1)K_2(T_2\mid x_2),
  \qquad
  U_{12}=U_1\otimes U_2.
\]
\end{proposition}

\begin{proof}
The product row has cardinality
\(d_{12}(x_1,x_2)=d_1(x_1)d_2(x_2)\), so the uniform coefficient factors.
The update factorization follows by separating the two finite sums.
\end{proof}

For an interacting control take a declared disjoint source sum with vertices
\(u_1,v_1\) in the first component and \(u_2,v_2\) in the second, and links
\[
  \partial_1e_1=[v_1]-[u_1],
  \qquad
  \partial_1e_2=[v_2]-[u_2].
\]
Let \(x_{uu}\) contain one \(\rho\)-occupation at each of \(u_1,u_2\), and let
\(x_{vv}\) contain one at each of \(v_1,v_2\). Both have sector
\((2,[0])\). The joint
transition
\[
  T_{12}=(x_{uu},e_1+e_2,0,\Pi_{\mathrm{joint}},x_{vv})
\]
is admissible because
\[
  q(x_{vv})-q(x_{uu})
  =
  [v_1]+[v_2]-[u_1]-[u_2]
  =
  \partial_1(e_1+e_2).
\]
Declare only
\[
  \Enabled_{x_{uu}}=\{\mathbf1_{x_{uu}},T_{12}\},
  \qquad
  \Enabled_{x_{vv}}=\{\mathbf1_{x_{vv}}\}.
\]
The projections of the first row contain an identity and a move in each
component. Their Cartesian product would have four records, including the two
mixed records, whereas the declared joint row contains only the two diagonal
records. The synchronized joint enablement and the omitted mixed records are
explicit source-control declarations, not an accidental finite truncation.
The row therefore does not factor as an independent product. This is a
candidate-level non-Cartesian interaction rule. It is not a physical force,
interaction, or coupling law.

\section{Null, degenerate, malformed, and underdetermined controls}

For the valid P7-T01 null configuration
\(x_0=(m=0,z=0)\), choose the sole row
\(\Enabled_{x_0}=\{\mathbf1_{x_0}\}\). Then
\[
  U\delta_{x_0}=\delta_{x_0}.
\]
This is a valid null control, not a physical vacuum or ground state.

For a valid non-null configuration \(x_d\), choose the sole row
\(\Enabled_{x_d}=\{\mathbf1_{x_d}\}\). Then
\(U\delta_{x_d}=\delta_{x_d}\). This is the degenerate identity-only branch.
It is distinct from the null configuration and does not establish physical
stasis.

Malformed candidate data return \(\Malformed\) with the failed condition:
\begin{enumerate}[leftmargin=*]
\item a purported \(x_u\to x_v\) transition with \(J=0\) has residual
  \([v]-[u]\) and fails P7-T01 continuity;
\item an empty row or a row omitting \(\mathbf1_x\) fails the totality
  convention;
\item negative coefficients or a row whose coefficients do not sum to one
  fail the kernel rule;
\item a support decision, coefficient, or cutoff defined from target metric
  distance, detector output, desired GR behavior, validator status, registry
  metadata, generated derivatives, or commit state fails the source-only
  guard.
\end{enumerate}
Malformed records never enter the kernel support.

If P7-T01 data are well typed but no finite control domain or enabled-row
declaration is supplied for a requested update, the result is
\(\Undetermined\), not zero, nullity, degeneracy, or malformedness. This
branch records exactly why current ontology does not derive a canonical
transition kernel.

\section{Symmetry, constraints, and conservation language}

The exact structure has three different sources:
\begin{longtable}{p{0.27\textwidth}p{0.31\textwidth}p{0.34\textwidth}}
\toprule
Result & Mathematical source & Authorized reading\\
\midrule
Equal \(Q,D\) across support &
P7-T01 chain constraints &
Inherited candidate-sector admissibility only\\
Row normalization and \(\|\mu\|_1\) preservation &
Uniform finite kernel &
New bookkeeping identity, not physical probability conservation\\
Invariant-subspace preservation &
Declared presentation-equivariance theorem &
Candidate presentation constraint, not target covariance or Noether law\\
Independent factorization &
Cartesian enabled rows &
Candidate composition identity, not empirical independence\\
Interacting nonfactorization &
Declared non-Cartesian joint row &
Proposal-only structural interaction, not force or coupling\\
\bottomrule
\end{longtable}

This is not a physical conservation law, and no physical conservation theorem
is claimed. The P7-T02 burden is met by an
explicit transition law plus its exact normalization, closure, sector-support,
and symmetry-compatibility constraints.

\section{Source-extension classification and open burdens}

The enabled-row assignment, identity convention, uniform weighting, formal
composition index, and composition rules are new proposal-only primitives.
Current ontology does not derive them. The positive construction shows that a
conservative, target-atlas-free source extension is mathematically available;
it does not make the extension canonical, empirically adequate, unique, or
adopted.

The following burdens remain open:
\begin{enumerate}[leftmargin=*]
\item protected selection or adoption of a source transition law;
\item physical time, probability, rate, persistence, or dynamics;
\item operational clocks, rods, signals, detectors, and free-fall systems;
\item continuum or principal-symbol analysis and common propagation;
\item effective geometry and universality;
\item matter action, stress energy, and universal coupling;
\item Einstein equations and exact-GR benchmark recovery;
\item every protected ontology, source-law, and benchmark decision.
\end{enumerate}

\begin{nonconclusion}[Exact claim boundary]
The decisive Candidate Constructor result is the machine status
\[
\texttt{constructed\_candidate}.
\]
The constructed object is exactly the
proposal-only finite kernel and update defined above. It is not a
canonical-ontology candidate, adopted source law, physical matter dynamics,
physical stochastic process, time evolution, physical conservation theorem,
detector model, metric theory, matter coupling, or GR derivation. Validator
and checkpoint success are operational evidence only.
\end{nonconclusion}

\end{document}
